\bookchapter{الفصل 31: مقدمة في تصوير الفيديوهات القصيرة وإنشاء المحتوى}

\chapterintro{في 2026، الفيديو القصير لم يعد “لعبة الشباب”، بل أصبح الطريقة الأساسية للجميع للحصول على المعلومات والتعبير عن الذات. سواء أردت توثيق الحياة أو مشاركة المعرفة أو عملاً جانبياً، هاتفك أداة الإبداع الأفضل — لا تحتاج كاميرا ولا تعلم Premiere، هاتف واحد يكفي للبدء. هذا الفصل يبدأ بك من الصفر.}

\booksection{واحد: فكر قبل التصوير: لمن تصور؟}

أول خطأ لكثيرين “يلتقط الهاتف ويصور”، بعد التصوير لا يعرف لمن. قبل البدء، أجب عن ثلاثة أسئلة:

\begin{enumerate}
  \item \textbf{من جمهوري؟}— التصوير للأطفال، للأقران، لكبار السن، أسلوب المحتوى مختلف تماماً.
  \item \textbf{ما القيمة التي أقدمها؟}— إضحاك (ترفيه)، تعليم (معرفة)، إثارة شهية (طعام)، إسعاد (عاطفة). واحدة على الأقل.
  \item \textbf{في أي منصة يمكنني التحديث المستمر؟}— Douyin/Kuaishou (ترفيه عام)، Xiaohongshu (نمط حياة)، Bilibili (محتوى عميق)، Video Account (دائرة اجتماعية). اختر واحدة كقاعدة رئيسية، لا توزع على كل المنصات من البداية.
\end{enumerate}

\booksection{اثنان: إعدادات التصوير الأساسية — 5 أمور قبل البدء}

\begin{enumerate}
  \item \textbf{الدقة 1080P، معدل الإطار 30fps}: هذا معيار منصات الفيديو القصير. 4K ملف كبير ورفع بطيء، 60fps لا ميزة على المنصات الاجتماعية وملف أكبر.
  \item \textbf{نظّف العدسة}: الأول والأهم. قبل التصوير امسح العدسة بثوبك، الصورة أوضح بمستوى.
  \item \textbf{قفل التركيز والتعريض}: في إطار العرض اضغط مطولاً على الهدف 2 ثانية، يظهر "AE/AF Lock". هكذا لا تتغير الصورة فاتحة وداكنة عند الحركة.
  \item \textbf{تصوير أفقي أم عمودي؟}— للنشر على Douyin/Kuaishou/Video Account/Xiaohongshu\textbf{صوّر عمودياً}(9:16)، للنشر على Bilibili/YouTube\textbf{صوّر أفقياً}(16:9). فكر في المنصة قبل تحديد الاتجاه.
  \item \textbf{فعّل خطوط الشبكة}: الإعدادات الكاميرا خطوط الشبكة. أساس تكوين قاعدة الثلث، يساعدك في إيجاد مركز الصورة بسرعة.
\end{enumerate}

\booksection{ثلاثة: خمس حركات كاميرا — اجعل الصورة “تتحرك”}

أكبر مشكلة للمبتدئين الصورة “ثابتة” — واقف طوال الوقت. تعلم خمس حركات أساسية، الصورة تصبح احترافية فوراً:

\begin{itemize}
  \item \textbf{دفع (Push)}: الهاتف يقترب ببطء من الهدف. يخلق توتراً أو حميمية "اقتراب".
  \item \textbf{سحب (Pull)}: الهاتف يبتعد ببطء عن الهدف. من التفصيل للبانوراما، يوضح البيئة.
  \item \textbf{تدوير (Pan)}: الهاتف ثابت، يدور يميناً ويساراً. مناسب للمناظر وعرض المساحات.
  \item \textbf{تحريك (Track)}: الهاتف يتحرك موازياً، يتبع الهدف. يحتاج ثبات اليد، يُفضل مع مثبت هاتف.
  \item \textbf{متابعة (Follow)}: الهاتف خلف أو أمام الهدف. حركة الكاميرا الأكثر استخداماً في Vlog.
\end{itemize}

\textbf{الثبات الأولوية الأولى}: الاهتزاز يبدو غير احترافي. الإمساك بكلتا اليدين أثبت من واحدة؛ عند المشي اثنِ الركبتين قليلاً، استخدم “خطوة النينجا” لتقليل الاهتزاز؛ إذا كانت الميزانية متاحة، اشترِ مثبت هاتف (سلسلة DJI OM حوالي 500-800 يوان)، تحول نوعي.

\booksection{أربعة: مونتاج للمبتدئين — JianYing يكفي}

JianYing (CapCut نسخة الصين) أنسب أداة مونتاج للمبتدئين على الهاتف، مجانية وكاملة والقوالب كثيرة. سير العمل الأساسي:

\begin{enumerate}
  \item \textbf{استيراد المواد}: افتح JianYing ابدأ الإبداع اختر مقاطع الفيديو المصوّرة واستوردها.
  \item \textbf{تقصير}: اسحب الخط الزمني، قص اللقطات غير المفيدة في بداية ونهاية كل مقطع. المبدأ: احتفظ فقط بالأجزاء ذات المعلومات، الأقصر أفضل من الأطول.
  \item \textbf{إضافة انتقال}: بين مقطعين اضغط زر "انتقال"، اختر انتقالاً بسيطاً. يُنصح للمبتدئين بـ "تداخل" أو "شاشة سوداء"، لا انتقالات ثلاثية الأبعاد استعراضية.
  \item \textbf{إضافة ترجمة}: JianYing لديه وظيفة "تمييز الترجمة"، يولّد الترجمة تلقائياً بضغطة. الدقة حوالي 90\%،راجع يدوياً.
  \item \textbf{إضافة موسيقى}: مكتبة موسيقى JianYing فيها كميات كبيرة من موسيقى خالية من حقوق النشر. اختر BGM يتناسب مع الإيقاع والمحتوى، اضبط الصوت لـ 1/3 صوت الفيديو الأصلي.
  \item \textbf{تصدير}: تصدير 1080P، معدل إطار 30fps، معدل البت "موصى به".
\end{enumerate}

\booksection{خمسة: أهمية الغلاف — 70\% يحكمون بالغلاف}

غلاف الفيديو القصير هو “الواجهة”، يحدد إن كانوا سيفتحونه. نقاط:

\begin{itemize}
  \item \textbf{الغلاف يجب أن يحوي جملة}: اكتب محتوى الفيديو الأساسي بخط كبير، مثل "3 دقائق لتعلم تحرير الصور بالهاتف". الخط كبير، اللون زاهٍ، تباين مع الخلفية.
  \item \textbf{الغلاف والمحتوى متناسقان}: لا تكن "صياد عناوين" — الغلاف يقول شيئاً، والداخل آخر، سيغادر المشاهد فوراً، والخوارزمية تخفض الترتيب.
  \item \textbf{تعبيرات الوجه حية}: إذا كان في الغلاف شخص، اختر لحظة تعبير حي أو مبالغ. الوجه يجذب النقر أكثر من النص.
  \item \textbf{صنع الغلاف في JianYing}: اختر أفضل إطار في JianYing، أضف نصاً متداخلاً، صدّر كصورة كغلاف.
\end{itemize}

\booksection{ستة: ثلاث قواعد لإنشاء المحتوى}

\begin{itemize}
  \item \textbf{أول 3 ثوانٍ تحدد المصير}: على منصات الفيديو القصير، يقرر المشاهد خلال 3 ثوانٍ المغادرة أو البقاء. أول 3 ثوانٍ يجب أن تقدم "خطافاً" — سؤالاً، تشويقاً، مشهداً غير متوقع.
  \item \textbf{فيديو واحد موضوع واحد}: خطأ المبتدئين حشو الكثير. فيديو واحد يحل مشكلة واحدة، الشرح العميق أهم من الشرح الشامل.
  \item \textbf{الاستمرارية أهم من الموهبة}: توصية الخوارزمية تحتاج "تغذية الحساب". أول 20 فيديو تقريباً بلا مشاهدات، هذا طبيعي. استمر بـ 2-3 فيديو أسبوعياً، ليتعرف الخوارزمية على اتجاه محتواك، والزيارات تزداد تدريجياً.
\end{itemize}

\begin{tcbpractice}{جرّب: صوّر فيديو قصير 30 ثانية}
\begin{enumerate}
  \item اختر معرفة صغيرة تتقنها (مثل "كيف تكبّر خط الهاتف").
  \item صوّر 3-5 مقاطع عمودياً بالهاتف، كل مقطع 5-10 ثوانٍ.
  \item استخدم JianYing للتوصيل والقص وإضافة الترجمة وBGM.
  \item اصنع غلافاً بنص كبير.
  \item بعد التصدير أرسله لصديق أولاً، اجمع الملاحظات ثم انشر.
\end{enumerate}

\end{tcbpractice}

\begin{tcbsummary}{خلاصة الفصل}
\begin{itemize}
  \item فكر قبل التصوير: من الجمهور، ما القيمة، في أي منصة
  \item 5 أمور قبل البدء: 1080P/30fps، نظف العدسة، قفل التركيز والتعريض، حدد الأفقي/العمودي، فعّل خطوط الشبكة
  \item خمس حركات كاميرا: دفع، سحب، تدوير، تحريك، متابعة — الثبات أولاً
  \item المونتاج بـ JianYing: قص انتقال ترجمة BGM تصدير، ست خطوات
  \item أول 3 ثوانٍ تحدد المصير، فيديو واحد موضوع واحد، الاستمرارية أهم من الموهبة
\end{itemize}

\end{tcbsummary}
